%% parts/p03/preservation-abstract-logic.tex
%% Cited survey results are distinguished from the complete Robinson, Beth,
%% and Lindstrom proofs routed through the Part III annex.

\chapter{Preservation and Abstract Model Theory}
\label{sec:p03-preservation-abstract}

\noindent\textit{Reader note.}  This is a separately authored survey of
preservation, joint consistency, and abstract-logical comparison after the
exact inherited Part~III sequence.  Its statements follow
\cite{ChangKeisler1990,Lindstrom1969}; they do not replace any inherited proof
or categoricity result.

This chapter collects three statement families whose common subject is the
transport of first-order information: between theories in overlapping
languages, along maps and chains of structures, and across extensions of the
first-order formalism.  The distinctions are syntactic and semantic, so the
hypotheses are part of each lookup entry rather than deferred to its citation.

\section{Existential closure and model companions}
\label{sec:p03-model-companions}

\begin{hbdefinition}{def:p03-existentially-closed}{Existentially closed model}
Let \(T\) be a theory.  A model \(M\entails T\) is
\emph{existentially closed among models of \(T\)} when, for every extension
\(M\subseteq N\entails T\), every existential \(\Lang(M)\)-formula true in
\(N\) is already true in \(M\).  Here an \emph{existential formula} has the
form \(\exists\bar y\,\theta(\bar x,\bar y)\) with \(\theta\)
quantifier-free; the tuple \(\bar y\) may be empty.  The phrase is needed to distinguish
existence of witnesses inside \(M\) from mere consistency of those witnesses
in an extension.  It is a property of a model relative to a specified theory
and embedding notion, not a synonym for algebraic closure.  The distinction
is used to identify the models of the companions below.
\index{existentially closed model}
\end{hbdefinition}

\begin{hbdefinition}{def:p03-model-companion}{Companions, model companions, and model completions}
Two theories in the same signature are \emph{companions} when they have the
same universal consequences.  A \emph{model companion} \(T^*\) of \(T\) is
a companion of \(T\) that is model-complete.  It is a
\emph{model completion} of \(T\) when, in addition, for every
\(M\entails T\), the theory
\(T^*\cup\operatorname{Diag}(M)\) is complete in the parameter language
\(\Lang(M)\).

These three names separate universal compatibility, elementarity of
embeddings between companion models, and uniqueness over a fixed base model.
They name relations between theories, not constructions of one preferred
model.  A model completion is therefore a model companion, but the definitions
are not interchangeable.  The stronger distinction is needed for the field
and ordered-field examples.
\index{model companion}\index{model completion}
\end{hbdefinition}

\begin{hbtheorem}{thm:p03-robinson-model-completeness}{cited=Robinson1956}{Robinson's test for model completeness}
For a theory \(T\), the following are equivalent:
\begin{enumerate}
  \item \(T\) is model-complete;
  \item whenever \(M\subseteq N\) are models of \(T\), the model \(M\) is
        existentially closed in \(N\); and
  \item every first-order formula is equivalent modulo \(T\) to an
        existential formula.
\end{enumerate}
\end{hbtheorem}

Model completeness immediately gives existential closure.  For the converse,
Robinson's diagram lemma says that if clause~2 holds for every pair of nested
\(T\)-models, then every formula is equivalent modulo \(T\) to an existential
formula: otherwise compactness applied to a renamed diagram produces nested
models that agree on all existential \(\Lang(M)\)-formulas but disagree on the
chosen formula.  This is the normal-form step; existential closure alone is
not a direct invocation of Tarski--Vaught.  Finally, clause~3 applied to both
\(\varphi\) and \(\neg\varphi\) makes nested \(T\)-models agree on every
formula, because existential formulas are preserved upward.  This proves the
test rather than using ``model-complete'' as a name for quantifier elimination
\cite{Robinson1956,ChangKeisler1990}.

\begin{hbtheorem}{thm:p03-field-model-completions}{cited=Marker2002}{Fields and ordered fields}
For each characteristic \(p\), where \(p=0\) or \(p\) is prime,
\(\mathsf{ACF}_p\) is the model completion of the theory of fields of
characteristic \(p\) in the ring language.  The theory \(\mathsf{RCF}\) is
the model completion of ordered fields in the ordered-ring language.
Consequently the existentially closed fields of characteristic \(p\) are the
algebraically closed fields, and the existentially closed ordered fields are
the real closed fields.
\end{hbtheorem}

If a field is not algebraically closed, a root of a nonconstant irreducible
polynomial appears in a field extension and witnesses failure of existential
closure.  An algebraic closure supplies an extension of every field, and
quantifier elimination for \(\mathsf{ACF}_p\) gives model completeness and
completeness over each atomic diagram.  The ordered-field argument uses a real
closure and quantifier elimination in the language that includes the order.
This also explains why the language columns in the quantifier-elimination
catalog are part of the theorem, not typography.

\section{Robinson joint consistency}
\label{sec:p03-robinson-joint-consistency}

Let \(\Lang_1\) and \(\Lang_2\) be first-order signatures and put
\(\Lang_0=\Lang_1\cap\Lang_2\).  Assume that every symbol in the overlap has
the same function/relation kind and the same arity in both signatures;
\(\Lang_0\) is the common restriction with those arities, and
\(\Lang_1\cup\Lang_2\) uses the compatible unions of the symbol sets and arity
maps.  For \(i\in\{1,2\}\), let \(T_i\) be a
consistent \(\Lang_i\)-theory.  The point of the following hypothesis is
agreement in the common language at the level of consequences; equality of
the literal sets \(T_1\cap T_2\) would not express it.

\begin{hblemma}{thm:p03-elementary-reduct-extension}{proved=annex:p3.abstract.4}{Extending a reduct together with its expansion}
Let \(\Lang_0\subseteq\Lang\), let \(M\) be a \(\Lang\)-structure, and
let \(N\) be a \(\Lang_0\)-structure with
\(M\mathbin{\upharpoonright}\Lang_0\preccurlyeq N\).
There is a \(\Lang\)-structure \(M^*\) such that
\[
 M\preccurlyeq M^*,
 \qquad
 N\preccurlyeq M^*\mathbin{\upharpoonright}\Lang_0,
\]
with both maps the inclusions on their underlying domains.
The construction uses elementary diagrams and first-order compactness
under the ambient Choice assumption.
\end{hblemma}

Elementarity of the reduct inclusion is essential.  In the unary
predicate language, let \(A\) have one element and interpret \(P\)
as empty; let \(B\) add a new element satisfying \(P\).
Then \(A\) is a substructure of \(B\), but
\(\operatorname{ElDiag}(A)\) contains \(\forall x\,\neg P(x)\),
whereas \(\operatorname{ElDiag}(B)\) contains \(P(c_b)\) for the new
element.  The two elementary diagrams are inconsistent together.

\begin{hbtheorem}{thm:p03-robinson-joint-consistency}{proved=annex:p3.abstract.1}{Robinson joint consistency theorem}
Under the ambient Choice assumption, if there is no
\(\Lang_0\)-sentence \(\sigma\) such that
\[
 T_1\derives\sigma
 \qquad\text{and}\qquad
 T_2\derives\neg\sigma,
\]
then \(T_1\cup T_2\), regarded as a theory in
\(\Lang_1\cup\Lang_2\), is consistent.
\end{hbtheorem}

A standard equivalent presentation assumes that both \(T_1\) and \(T_2\)
extend one complete \(\Lang_0\)-theory.  This is a joint-consistency statement
about theories in a fixed common signature, not an assertion that two given
models can be amalgamated over a specified common substructure
\cite{RobinsonJointConsistency1956}.

The proof first obtains a complete common theory compatible with both
sides.  Elementary diagrams place a model of one side inside the reduct
of a model of the other.  The preceding lemma then alternates the two
extensions:
\[
 M_0\mathbin{\upharpoonright}\Lang_0\preccurlyeq N_0\mathbin{\upharpoonright}\Lang_0
 \preccurlyeq M_1\mathbin{\upharpoonright}\Lang_0\preccurlyeq N_1\mathbin{\upharpoonright}\Lang_0
 \preccurlyeq\cdots .
\]
The full \(\Lang_1\)-structures form one elementary chain and the full
\(\Lang_2\)-structures form another.  Their domains have the same union,
and their interpretations of every common symbol agree there.
The elementary chain theorem supplies the required joint model.
Annex~\ref{annex:p3.abstract.1} also records the alternative route
through isomorphic ultrapowers.

For example, let
\[
 T_1=\{\exists x\,P(x),\ \exists x\,\neg P(x)\},
 \qquad T_2=\{\forall x\,y\,(x=y)\}.
\]
Both theories are consistent, and their literal intersection is empty.
Their union is inconsistent: \(T_1\) requires two distinct elements,
whereas \(T_2\) requires exactly one.
The common signature has no nonlogical symbol, but the equality sentence
\(\exists x\,y\,(x\ne y)\) separates the two theories.
Logical equality therefore remains part of the common-language test.

\section{Preservation under structural change}
\label{sec:p03-preservation}

The preservation theorems below are classical results over \emph{all}
nonempty first-order structures.  Their finite-structure analogues can fail.

\begin{hbdefinition}{def:p03-elementary-class}{Elementary class}
A class \(\mathcal K\) of \(\Lang\)-structures is \emph{elementary} iff there
is a set \(T\) of \(\Lang\)-sentences such that
\(\mathcal K=\{\mathcal M:\mathcal M\entails T\}\).  It is
\emph{axiomatizable by} a syntactic fragment \(F\) iff such a \(T\) can be
chosen from \(F\).  This class-level notion describes a collection of whole
structures; it is distinct from the relation
\(\mathcal M\preccurlyeq\mathcal N\) between two structures.
\end{hbdefinition}

\begin{hbdefinition}{def:p03-preservation-relations}{Preservation under substructures, extensions, and chain unions}
A class \(\mathcal K\) is \emph{closed under substructures} iff
\(\mathcal N\in\mathcal K\) and \(\mathcal M\subseteq\mathcal N\) imply
\(\mathcal M\in\mathcal K\).  A sentence \(\varphi\) is
\emph{preserved under extensions} iff
\(\mathcal M\subseteq\mathcal N\) and
\(\mathcal M\entails\varphi\) imply
\(\mathcal N\entails\varphi\).  Finally, \(\mathcal K\) is
\emph{closed under unions of chains} iff the compatible union of every
nonempty ordinal-indexed chain of members of \(\mathcal K\), ordered by
substructure inclusion, again belongs to \(\mathcal K\).  This last condition
uses arbitrary substructure chains, not only elementary chains.
\index{preservation theorem}
\index{union of chains}
\end{hbdefinition}

\begin{hbdefinition}{def:p03-positive-formula}{Positive and existential-positive formulas}
In a relational signature with equality, a \emph{positive formula} is built
from \(\bottom\) and atomic formulas, including equality atoms, using only
\(\land,\lor,\forall,\exists\).  An \emph{existential-positive formula} is
built from \(\bottom\) and atoms using only \(\land,\lor,\exists\).  Thus
positive does not mean
existential-positive.  A formula \(\varphi(\bar x)\) is
\emph{preserved under surjective homomorphisms} iff, for every surjective
homomorphism \(h:\mathcal M\to\mathcal N\) and tuple \(\bar a\) from
\(\mathcal M\),
\[
 \mathcal M\entails\varphi(\bar a)
 \quad\Longrightarrow\quad
 \mathcal N\entails\varphi(h\bar a).
\]
\index{positive formula}
\index{existential-positive formula}
\end{hbdefinition}

\begin{hbtheorem}{thm:p03-los-tarski}{cited=ChangKeisler1990}{Los--Tarski preservation theorem@\L o\'s--Tarski preservation theorem}
For an elementary class \(\mathcal K\) of \(\Lang\)-structures, the following
are equivalent:
\begin{enumerate}
 \item \(\mathcal K\) is closed under substructures;
 \item \(\mathcal K\) is axiomatizable by universal \(\Lang\)-sentences.
\end{enumerate}
Equivalently at sentence level, a first-order sentence is preserved under
extensions iff it is logically equivalent, over all \(\Lang\)-structures, to
an existential sentence.
\end{hbtheorem}

\begin{hbtheorem}{thm:p03-chang-los-suszko}{cited=ChangKeisler1990}{Chang--Los--Suszko theorem@Chang--\L o\'s--Suszko theorem}
For an elementary class \(\mathcal K\) of \(\Lang\)-structures, the following
are equivalent:
\begin{enumerate}
 \item \(\mathcal K\) is closed under unions of chains of substructures;
 \item \(\mathcal K\) is axiomatizable by sentences of the form
       \(\forall\bar x\,\exists\bar y\,\theta(\bar x,\bar y)\), where
       \(\theta\) is quantifier-free and either quantifier block may be empty.
\end{enumerate}
\end{hbtheorem}

\begin{hbtheorem}{thm:p03-lyndon-preservation}{cited=Lyndon1959}{Lyndon positivity theorem}
In a relational signature with equality, a first-order formula is preserved
under surjective homomorphisms between all structures iff it is logically
equivalent to a positive formula.
\end{hbtheorem}

Table~\ref{tab:p03-preservation} keeps the syntactic fragments paired with the
semantic operation they characterize.  In particular, Lyndon's surjective
homomorphism theorem is not the distinct theorem pairing arbitrary
homomorphisms with existential-positive formulas
\cite{ChangKeisler1990,Lyndon1959}.

\begingroup\small
\begin{longtable}{@{}p{0.21\textwidth}p{0.31\textwidth}p{0.35\textwidth}@{}}
\caption{Preservation theorems over all structures.}
\label{tab:p03-preservation}\\
\toprule
Result & Semantic condition & Syntactic characterization \\
\midrule
\endfirsthead
\toprule
Result & Semantic condition & Syntactic characterization \\
\midrule
\endhead
{\raggedright \L o\'s--Tarski, class form\par} &
  {\raggedright Elementary class closed under substructures\par} &
  {\raggedright Axiomatizable by universal sentences\par} \\
\addlinespace
{\raggedright \L o\'s--Tarski, sentence dual\par} &
  {\raggedright Sentence preserved under extensions\par} &
  {\raggedright Equivalent to an existential sentence\par} \\
\addlinespace
{\raggedright Chang--\L o\'s--Suszko\par} &
  {\raggedright Elementary class closed under unions of substructure chains\par} &
  {\raggedright Axiomatizable by \(\forall^{*}\exists^{*}\) sentences with
  quantifier-free matrix\par} \\
\addlinespace
{\raggedright Lyndon positivity\par} &
  {\raggedright Formula preserved under surjective homomorphisms in a
  relational signature with equality\par} &
  {\raggedright Equivalent to a positive formula using
  \(\land,\lor,\forall,\exists\)\par} \\
\bottomrule
\end{longtable}
\endgroup

The restriction to all structures is load-bearing: in particular, the
substructure form of \L o\'s--Tarski is not valid when the universe of discourse
is restricted to finite structures.

\section{Beth definability and Craig interpolation}
\label{sec:p03-beth-craig}

\begin{hbdefinition}{def:p03-implicit-explicit-definability}{Implicit and explicit definability}
Let \(P\) be a new relation symbol of finite arity \(n\geq0\), and let \(T\) be a theory in
\(\Lang\cup\{P\}\).  The theory \emph{implicitly defines \(P\) over
\(\Lang\)} when any two models of \(T\) with the same \(\Lang\)-reduct have
the same interpretation of \(P\).  It \emph{explicitly defines \(P\)} when
there is a \(\Lang\)-formula \(\theta(\bar x)\) such that
\[
 T\entails\forall\bar x\,
   \bigl(P(\bar x)\leftrightarrow\theta(\bar x)\bigr).
\]
Implicit definability is a uniqueness condition on expansions; explicit
definability supplies an eliminable formula in the old language.  Both are
relations between a theory, a symbol, and a reduct language, not new semantic
consequence relations.  Their nonredundancy is precisely the question settled
by Beth's theorem, and the distinction is used in its proof.
\index{Beth definability}\index{implicit definability}
\end{hbdefinition}

\begin{hbtheorem}{thm:p03-beth-definability}{proved=annex:p3.abstract.2}{Beth definability theorem}
Under the ambient Choice assumption, every relation symbol implicitly
defined by a first-order theory is explicitly definable over the same
reduct language.  The theory and signature may be infinite.
One finite \(T_0\subseteq T\) already proves an explicit definition
\[
 \forall\bar x\,
       (P(\bar x)\leftrightarrow\theta(\bar x)),
 \qquad \FV(\theta)\subseteq\{x_1,\ldots,x_n\}.
\]
For \(n=0\), the defining formula is a sentence and the displayed
quantifier block is empty.
\end{hbtheorem}

\begin{hbtheorem}{thm:p03-craig-robinson-beth}{proved=annex:p3.abstract.5}{Craig--Robinson equivalence and Beth definability}
For classical first-order logic, using compactness, Craig interpolation
and the common-sentence form of Robinson joint consistency imply one
another.  Each implies Beth definability.
\end{hbtheorem}

The model-theoretic route starts with compactness, elementary diagrams,
and elementary chains, proves joint consistency, and obtains interpolation
by its contrapositive.  Renaming the implicitly defined predicate then gives
Beth's theorem.  This proof uses the same common-language and free-variable
conditions as Theorem~\ref{thm:craig-interpolation}, whose separate
Maehara proof is in Part~VII.

The language of an explicit definition matters.  Suppose \(A\models B\)
and \(q\) is a fresh nullary predicate.  In the union of the two
languages, the sentence
\[
 H(q)=(q\land A)\lor(\neg q\land\neg B)
\]
implicitly defines \(q\), with \(H(q)\models q\leftrightarrow A\).
If \(q\) is true, \(A\) follows from the first disjunct; if \(q\)
is false, \(\neg B\) and \(A\models B\) imply \(\neg A\).
This definition may use symbols private to \(A\).
For instance, if \(A=P\land Q\) and \(B=P\lor R\), the only shared
nonlogical symbol is \(P\), but \(A\) also contains \(Q\).
An explicit definition in the union language therefore does not
by itself supply an interpolant in the intersection language.

The historical sources are
\cite{Beth1953,Craig1957,RobinsonJointConsistency1956,Hodges1993}.

\section{Infinitary and generalized-quantifier landmarks}
\label{sec:p03-abstract-landmarks}

The syntax and satisfaction clauses for \(L_{\omega_1\omega}\) are fixed in
Definition~\ref{def:p02-lomega-one-omega}.  Its compactness failure, Scott
sentences for countable structures, and the countable-model restriction are
Theorem~\ref{thm:p02-lomega-one-omega-calibration}; the exact infinitary proof
boundary is Theorem~\ref{thm:p02-karp-completeness}.  Those public definitions
are not repeated here.  Their role in abstract model theory is diagnostic:
\(L_{\omega_1\omega}\) has strictly more syntax than first-order logic, but
loses compactness, while Scott sentences identify countable structures only
within the countable models \cite{Karp1964,Scott1965}.

Definition~\ref{def:p02-generalized-quantifier} treats a generalized
quantifier as isomorphism-invariant semantic data.  The most important
cardinality instance for this chapter is the following one.

\begin{hbdefinition}{def:p03-q-one}{The uncountability quantifier \(Q_1\)}
The unary generalized quantifier \(Q_1\) is the class of pairs \((M,A)\) with
\(M\ne\varnothing\), \(A\subseteq M\), and \(A\) uncountable.  Thus
\[
 M,s\entails Q_1x\,\varphi
 \quad\Longleftrightarrow\quad
 \{a\in M:M,s[x\mapsto a]\entails\varphi\}\text{ is uncountable}.
\]
The subscript distinguishes this fixed cardinality quantifier from an
arbitrary member of a generalized-quantifier family.  It is an
object-language quantifier with the displayed semantic parameter, not a new
predicate on \(M\), and it differs from ``there are infinitely many.''  The
distinction is required by Keisler's completeness theorem.
\index{uncountability quantifier}\index{Q1@\(Q_1\)}
\end{hbdefinition}

\begin{hbtheorem}{thm:p03-keisler-q-one}{cited=Keisler1970}{Keisler completeness for \(L(Q_1)\)}
There is a sound complete axiomatization of validity for first-order logic
extended by \(Q_1\).  In particular, every countable set of
\(L(Q_1)\)-sentences whose finite subsets are satisfiable has a model.  In a
countable signature every satisfiable \(L(Q_1)\)-sentence has a model of
cardinality at most \(\aleph_1\).
\end{hbtheorem}

The upper bound cannot be lowered to \(\aleph_0\), because
\(Q_1x\,(x=x)\) has no countable model.  Thus this proper extension retains a
countable compactness theorem but fails the countable downward
L\"owenheim--Skolem property used by Lindstr\"om.  The theorem concerns this
one quantifier and its stated semantics; it is not a completeness theorem for
all generalized quantifiers \cite{Mostowski1957,Lindstrom1966,Keisler1970}.

\section{Lindstr\"om's characterization}
\label{sec:p03-lindstrom}

The term \emph{regular abstract logic} has exactly the metalinguistic standing
fixed in the Part-2 first-order metatheorem summary
(Section~\ref{sec:p02-fol-metatheorem-summary}).  Explicitly, for each
relational signature \(\tau\), it assigns a class of sentences and an
associated satisfaction relation invariant under isomorphism.  Every sentence
has a finite occurrence signature, satisfaction depends only on the reduct to
that signature, and arity-preserving renamings transport both sentences and
satisfaction.  The logic contains relational first-order logic and is closed
under negation and conjunction.  Abstraction of a fresh constant has the usual
first-order quantifier semantics; for a fresh unary predicate \(U\), unary
relativization \(\varphi^U\) is evaluated in the induced \(\tau\)-structure on
the required nonempty set \(U^{\mathcal M}\).  This package is a condition on a
logic-and-satisfaction assignment, not an object-language connective or the
proof-theoretic fragment also called regular logic.  It is the Part-2 notion,
deliberately relocated rather than redefined here
\cite[Chapter~XIII, \S1]{EbbinghausFlumThomas1994}.

\begin{hbdefinition}{def:p03-potential-isomorphism}{Potential isomorphism}
Two structures \(A\) and \(B\) in the same relational signature are
\emph{potentially isomorphic} when there is a nonempty family \(\mathcal I\)
of finite partial isomorphisms from \(A\) to \(B\) such that, for every
\(p\in\mathcal I\), each element of \(A\) can be put into the domain of an
extension \(q\in\mathcal I\), and each element of \(B\) can be put into the
range of such an extension.  This relation is needed because agreement in
each fixed finite EF game does not itself supply one isomorphism.  It is a
metalinguistic relation between structures, not a map in the signature.  It
is weaker than isomorphism for arbitrary uncountable structures but, by
back-and-forth, equivalent to isomorphism for countable structures.  Those two
facts are used separately in the proof below.
\index{potential isomorphism}
\end{hbdefinition}

\begin{hbtheorem}{thm:p03-lindstrom}{proved=annex:p3.abstract.3}{Lindstrom's theorem@Lindstr\"om's theorem}
Let \(\mathcal L^{\ast}\) be a regular abstract logic in that fixed relational
sense.  Assume:
\begin{enumerate}
 \item whenever a countable set of \(\mathcal L^{\ast}\)-sentences has every
       finite subset satisfiable, the whole set is satisfiable; and
 \item every satisfiable \(\mathcal L^{\ast}\)-sentence has a countable model.
\end{enumerate}
Then every \(\mathcal L^{\ast}\)-sentence is equivalent, over its finite
occurrence signature, to a relational first-order sentence.  Conversely,
first-order sentences already belong to \(\mathcal L^{\ast}\).
\end{hbtheorem}

The annex supplies the finite sequence code, the potential-isomorphism lemma,
the bounded-rank lemma, and the nonstandard-index compactness construction.
It uses countable compactness and the sentence-level downward
L\"owenheim--Skolem hypothesis exactly as stated
\cite{Lindstrom1969,EbbinghausFlumThomas1994}.

\section{Independence-friendly and dependence logics}
\label{sec:p03-dependence-logics}

\begin{hbdefinition}{def:p03-team-semantics}{Teams and dependence semantics}
A \emph{team} \(X\) on a structure \(M\) is a set of assignments with one
common finite variable domain.  The dependence-logic grammar is generated by
first-order literals, \(\bottom\), positive atoms \(=(\bar x,y)\), and
\(\land,\lor,\exists,\forall\).  Negated dependence atoms are not formulas;
first-order negation occurs only inside literals.  Satisfaction is defined on
teams by
\[
\begin{aligned}
 M,X\entails\alpha
   &\Longleftrightarrow (\forall s\in X)\,M,s\entails\alpha
        &&\text{for first-order literals }\alpha,\\
 M,X\entails\bottom
   &\Longleftrightarrow X=\varnothing,\\
 M,X\entails\varphi\land\psi
   &\Longleftrightarrow M,X\entails\varphi\text{ and }M,X\entails\psi,\\
 M,X\entails\varphi\lor\psi
   &\Longleftrightarrow X=Y\cup Z\text{ for some }Y,Z
       \text{ with }M,Y\entails\varphi, M,Z\entails\psi,\\
 M,X\entails\exists x\,\varphi
   &\Longleftrightarrow M,X[F/x]\entails\varphi
       \text{ for some }F:X\to\mathcal P(M)\setminus\{\varnothing\},\\
 M,X\entails\forall x\,\varphi
   &\Longleftrightarrow M,X[M/x]\entails\varphi,\\
 M,X\entails=(\bar x,y)
   &\Longleftrightarrow
     (\forall s,s'\in X)\bigl(s(\bar x)=s'(\bar x)\to s(y)=s'(y)\bigr),
\end{aligned}
\]
where
\(X[F/x]=\{s[a/x]:s\in X, a\in F(s)\}\) and
\(X[M/x]=\{s[a/x]:s\in X, a\in M\}\).

Teams are needed because functional dependence is a relation among possible
assignments and cannot be witnessed by one assignment.  A team is a semantic
set of assignments, not a new model or object-language relation; the
dependence atom differs from an ordinary relation atom by comparing pairs of
assignments.  The distinction is used in the worked sentence and expressive-
power theorem below.
\index{dependence logic}\index{team semantics}
\end{hbdefinition}

\begin{hbdefinition}{def:p03-if-logic}{Independence-friendly quantification}
Independence-friendly logic permits a slashed existential quantifier
\((\exists x/W)\varphi\).  In the semantic evaluation game, the existential
player's strategy for choosing \(x\) must be uniform on positions that agree
outside the variables in \(W\); equivalently, the choice may not depend on the
values hidden by the slash set.  This quantifier is needed to state
independence constraints directly in the quantifier prefix.  It is a
game-semantic restriction on a move, not a new generalized cardinality
quantifier, and differs syntactically from the dependence atom even though the
sentence-level expressive powers coincide below.
\index{independence-friendly logic}
\end{hbdefinition}

For example,
\[
 \forall x\,\forall z\,\exists y\,
   \bigl(=(x,y)\land\varphi(x,z,y)\bigr)
\]
says that one function of \(x\) alone chooses \(y\), independently of \(z\),
so that \(\varphi\) holds for every pair \((x,z)\).  After the two universal
expansions there are assignments with the same \(x\) and different \(z\);
the dependence atom forces all of them to receive the same \(y\).  It is
therefore doing work that the pointwise first-order matrix cannot do.

\begin{hbtheorem}{thm:p03-dependence-eso}{cited=Vaananen2007}{Sentence-level expressive power}
Over structures, the properties definable by dependence-logic sentences are
exactly those definable by existential second-order sentences.  The same
sentence-level expressive class is obtained by independence-friendly logic
under its game semantics.
\end{hbtheorem}

The equivalence translates team choices into existentially quantified
relations or Skolem functions, and conversely translates those witnesses into
dependence atoms or uniform game strategies.  It is a sentence-level theorem;
open formulas retain team information and are not identified here with
ordinary Tarskian formulas \cite{HintikkaSandu1989,Vaananen2007}.
